% Options for packages loaded elsewhere
\PassOptionsToPackage{unicode}{hyperref}
\PassOptionsToPackage{hyphens}{url}
%
\documentclass[
  10pt,
  a4paper,
]{article}
\usepackage{amsmath,amssymb}
\usepackage{iftex}
\ifPDFTeX
  \usepackage[T1]{fontenc}
  \usepackage[utf8]{inputenc}
  \usepackage{textcomp} % provide euro and other symbols
\else % if luatex or xetex
  \usepackage{unicode-math} % this also loads fontspec
  \defaultfontfeatures{Scale=MatchLowercase}
  \defaultfontfeatures[\rmfamily]{Ligatures=TeX,Scale=1}
\fi
\usepackage{lmodern}
\ifPDFTeX\else
  % xetex/luatex font selection
  \setmainfont[]{DejaVu Serif}
  \setmonofont[]{DejaVu Sans Mono}
  \setmathfont[]{Latin Modern Math}
\fi
% Use upquote if available, for straight quotes in verbatim environments
\IfFileExists{upquote.sty}{\usepackage{upquote}}{}
\IfFileExists{microtype.sty}{% use microtype if available
  \usepackage[]{microtype}
  \UseMicrotypeSet[protrusion]{basicmath} % disable protrusion for tt fonts
}{}
\makeatletter
\@ifundefined{KOMAClassName}{% if non-KOMA class
  \IfFileExists{parskip.sty}{%
    \usepackage{parskip}
  }{% else
    \setlength{\parindent}{0pt}
    \setlength{\parskip}{6pt plus 2pt minus 1pt}}
}{% if KOMA class
  \KOMAoptions{parskip=half}}
\makeatother
\usepackage{xcolor}
\usepackage[margin=20mm]{geometry}
\setlength{\emergencystretch}{3em} % prevent overfull lines
\providecommand{\tightlist}{%
  \setlength{\itemsep}{0pt}\setlength{\parskip}{0pt}}
\setcounter{secnumdepth}{-\maxdimen} % remove section numbering
\ifLuaTeX
\usepackage[bidi=basic]{babel}
\else
\usepackage[bidi=default]{babel}
\fi
\babelprovide[main,import]{english}
\ifPDFTeX
\else
\babelfont{rm}[]{DejaVu Serif}
\fi
% get rid of language-specific shorthands (see #6817):
\let\LanguageShortHands\languageshorthands
\def\languageshorthands#1{}
\usepackage{amsmath,amssymb,amsthm,mathtools,graphicx,booktabs,longtable,needspace}
\usepackage{fvextra}
\DefineVerbatimEnvironment{Highlighting}{Verbatim}{commandchars=\\\{\},breaklines=true,breakanywhere=true,fontsize=\small}
\setlength{\emergencystretch}{3em}
\allowdisplaybreaks
\raggedbottom\sloppy
\widowpenalty=10000\clubpenalty=10000
\providecommand{\Tr}{\operatorname{Tr}}
\providecommand{\tr}{\operatorname{Tr}}
\ifLuaTeX
  \usepackage{selnolig}  % disable illegal ligatures
\fi
\usepackage{bookmark}
\IfFileExists{xurl.sty}{\usepackage{xurl}}{} % add URL line breaks if available
\urlstyle{same}
\hypersetup{
  pdftitle={Mathematical Terminology},
  pdfauthor={Kirtus Justus (Theoretical Physicist)  31 October 2026Port Coquitlam, British Columbia, Canada},
  pdflang={en},
  hidelinks,
  pdfcreator={LaTeX via pandoc}}

\title{Mathematical Terminology}
\author{Kirtus Justus (Theoretical Physicist) \textbullet{} 31 October 2026\tabularnewline Port Coquitlam, British Columbia, Canada}
\date{}

\begin{document}
\maketitle

This glossary fixes objects, conventions and limiting regimes. A named law or estimate is described here to identify what the term means; its proof status, hypotheses and quantitative conclusion belong to the cited paper and to the separate results guide. Numerical observations do not enter these definitions.

\subsection{Band-flattened Hamiltonian}\label{band-flattened-hamiltonian}

A self-adjoint gapped one-particle Hamiltonian replaced by its spectral sign or signed band projectors, preserving eigenprojectors but not the original excitation-energy magnitudes.

\nopagebreak[4]

\emph{Reference:} Projector-Equivalent Wilson--Dirac Hamiltonians and Dynamical Identification, §1.

\subsection{Boundary leakage}\label{boundary-leakage}

For a full-line self-adjoint covariance difference \(T\) and interval projection \(\chi\), the positive boundary operator \(\chi T(I-\chi)T\chi\) equals \(\chi T^2\chi-(\chi T\chi)^2\). Its trace is the squared Hilbert--Schmidt norm of the cross-boundary block.

\nopagebreak[4]

\emph{Reference:} Fourier-Space Boundary Defects for Wilson-Dirac Covariances, Theorem 1.

\subsection{Boundary-corrected longitudinal trace}\label{boundary-corrected-longitudinal-trace}

The longitudinal trace \(\operatorname{Tr}(S_NE_N)\) expressed using the interval compression of the full-line projector/majorant product and its cross-boundary terms. The boundary correction is needed because compression does not preserve products. The identity and error estimates belong to the cited theorem.

\nopagebreak[4]

\emph{Reference:} Boundary-Corrected Longitudinal Modular Reduction for a Wilson--Dirac Interval, §2.

\subsection{CAR algebra}\label{car-algebra}

The C*-algebra generated by creation and annihilation operators over a specified complex one-particle Hilbert space. Two annihilation operators anticommute; the anticommutator of an annihilation operator and a creation operator equals their one-particle inner product times the identity, in the declared convention.

\nopagebreak[4]

\emph{Reference:} Restricted Observables and CAR-State Distinguishability, §1.

\subsection{Coefficient-level edge matching}\label{coefficient-level-edge-matching}

The limiting relation between rescaled formal period coefficients and Stirling coefficients under a specified edge coordinate. It does not justify exchanging the bulk and edge limits without uniform remainders.

\nopagebreak[4]

\emph{Reference:} Higher-Order Modular Periods for the Half-Filled Sine Kernel, Edge poles and coefficient-level Stirling matching.

\subsection{Concentration occupation and modular sign}\label{concentration-occupation-and-modular-sign}

For a common eigenvector of \(B_N\) and \(J_N\), write \(\zeta\in(0,1)\) and \(\lambda\) for the eigenvalues. The conventions are \(\varepsilon=\log[(1-\zeta)/\zeta]\), \(t=\lambda/N\), and \(e=\varepsilon+N\eta(t)\). A particle--hole partner uses \((-\lambda,1-\zeta,-\varepsilon,-e)\).

\nopagebreak[4]

\emph{Reference:} Uniform Bulk Modular Asymptotics and the Leading Wilson Entropy Law, Statement and conventions.

\subsection{Conditional entropy finite-part rate}\label{conditional-entropy-finite-part-rate}

A convergence-rate conclusion for the remainder after subtracting both the leading logarithm and a proposed finite constant from the fixed-scaled-mass relative entropy. The word conditional refers to explicitly listed uniform comparison and convergence hypotheses. It does not weaken or strengthen a separate theorem about the leading logarithm.

\nopagebreak[4]

\emph{Reference:} Relative Entropy of Finite-Interval Wilson--Dirac Quasifree States, Conditional expansion and remainder analysis.

\subsection{Conformal majorant boundary energy}\label{conformal-majorant-boundary-energy}

The homogeneous half-Sobolev Fourier energy \(4\sum_{d\ge1}d|c_d(r)|^2\) for the scalarized Wilson-majorant coefficients \(c_d(r)\). Its conformal invariance and value are separate results; the quantity itself is a boundary energy rather than a modular trace.

\nopagebreak[4]

\emph{Reference:} Conformal Invariance and Sharp Wilson Majorant Boundary Bounds, §2.

\subsection{Continuum Wilson-majorant boundary leakage}\label{continuum-wilson-majorant-boundary-leakage}

The squared Hilbert--Schmidt norm of the positive Wilson majorant block crossing the interval boundary. The continuum quantity is defined by \(\Phi(\mu)=\frac{2}{\pi^2}\int_0^\infty\min(\mu,2x)[\int_0^\infty(1-t/\sqrt{1+t^2})\cos(2xt)\,dt]^2dx\) for \(\mu>0\). Its identification as a limit, positivity, monotonicity and concavity are results in the cited source.

\nopagebreak[4]

\emph{Reference:} Continuum Limit of the Wilson Majorant Boundary Defect, Definitions and Theorem 1.

\subsection{Covariance (quasi-free CAR)}\label{covariance-quasi-free-car}

A positive contraction \(0\le C\le I\) on the one-particle space that encodes the normal two-point functions of a gauge-invariant quasifree CAR state in a fixed convention. General fermionic quasifree states may also have anomalous pairing correlations. Two index conventions occur in this collection. If \(C^{\rm occ}_{ij}=\omega(a_i^\dagger a_j)\) and \(C^{\rm op}_{ij}=\omega(a_j^\dagger a_i)\), then \(C^{\rm op}=(C^{\rm occ})^{\mathsf T}=\overline{C^{\rm occ}}\). Operator identities and one-particle block reconstruction use the stated local convention. Determinants, spectra and \(|C_{12}|^2\) are unchanged by this translation. A change of index convention is not an experimental basis rotation or a physical change of state.

\nopagebreak[4]

\emph{Reference:} Operational Identification and Calibration in Two-Mode Quasifree Systems, §1; Restricted Observables and CAR-State Distinguishability, ``A covariance does not determine every gauge-invariant state''.

\subsection{Cross-boundary perturbation}\label{cross-boundary-perturbation}

The massive-minus-massless difference of the projection block mapping the outside of the unit interval into the inside. Its trace norm is controlled after the lattice and continuum operators have been embedded in the same Hilbert space. It is not the full covariance difference.

\nopagebreak[4]

\emph{Reference:} Cross-Boundary Trace-Class Estimates for Wilson--Dirac Quasifree States, Model and definitions.

\subsection{Detector-frame alias}\label{detector-frame-alias}

Two distinct parameter-and-detector-orientation settings producing the same observation probabilities under an explicitly declared measurement class. Injectivity may be restored by restricting and calibrating detector orientations; an alias for a larger class does not invalidate an inverse theorem proved on the smaller class.

\nopagebreak[4]

\emph{Reference:} Operational Identification and Calibration in Two-Mode Quasifree Systems, §§2 and 6.

\subsection{Dispersion deformation}\label{dispersion-deformation}

A specified change to momentum-dependent energy eigenvalues. Multiplication by a strictly positive scalar function preserves spectral projectors while it can change transition frequencies and energy ratios.

\nopagebreak[4]

\emph{Reference:} Projector-Equivalent Wilson--Dirac Hamiltonians and Dynamical Identification, §1.

\subsection{Effective continuum mass}\label{effective-continuum-mass}

The branch-aligned lattice parameter \(a_L=L\log(1+\mu/L)\). Its use in the cross-boundary kernel comparison precedes a separate estimate controlling the replacement of \(a_L\) by \(\mu\).

\nopagebreak[4]

\emph{Reference:} Cross-Boundary Trace-Class Estimates for Wilson--Dirac Quasifree States, Both boundary channels and direct continuum identification.

\subsection{Endpoint-normalized connection coefficient}\label{endpoint-normalized-connection-coefficient}

For \(0<\theta<\pi\), set \(p_\theta(x)=\cos x-\cos\theta\) and \(\alpha=(N^2-1)/4\). In \((p_\theta u^{\prime})^{\prime}+(\alpha\cos x-b)u=0\) on \((-\theta,\theta)\), \(C_\theta(b)\) denotes the coefficient of the right-endpoint logarithmic solution when both regular endpoint solutions are normalized to one. A prime on \(C_\theta\) means differentiation in \(b\). At a common finite concentration/Jacobi mode, \(\beta\) is the centered Fourier differential parameter; at half filling \(\beta=N\lambda/2\). The complementary-band identity uses this normalization.

\nopagebreak[4]

\emph{Reference:} Concentration Eigenvalues from Logarithmic Connection Coefficients, A normalized connection coefficient and its derivative.

\subsection{Exact spectral-tail concentration function}\label{exact-spectral-tail-concentration-function}

For the positive omitted differential eigenvalues \(\mu_{N,k}\), define \(F_N(b)=2\sum_{k\ge N}\operatorname{artanh}(b/\mu_{N,k})\) on \(|b|<\mu_{N,N}\). Convergence and its equality with a finite concentration logit are established in the cited spectral-tail theorem. The expansion in \(b\) and the large-\(N\) period expansion are distinct.

\nopagebreak[4]

\emph{Reference:} Spectral Tails and Exact Concentration Logits, Spectral-tail formula.

\subsection{Extended-Chebyshev Hermite dual}\label{extended-chebyshev-hermite-dual}

A linear combination of observation functions with prescribed simple and double contact zeros, and a proved sign on the allowed support half-line. When paired with a positive feasible representing measure and the appropriate error-set support condition, it can prove a global extremum. A local numerical root without this sign proof is insufficient.

\nopagebreak[4]

\emph{Reference:} Positive Spectral Measures Under Finite Moment and Laplace Constraints, §2.

\subsection{Extreme reference-modular discrepancy}\label{extreme-reference-modular-discrepancy}

The scalar discrepancy \(e_N=\log[(1-b_N)/b_N]+N\eta(\lambda_N/N)\) at the reference mode paired with the largest Jacobi eigenvalue \(\lambda_N\). A fixed-index asymptotic for this mode does not define or control the full discrepancy operator norm.

\nopagebreak[4]

\emph{Reference:} Half-Filled Sine-Kernel Endpoint and Extreme Modular Discrepancy, §§1 and 5.

\subsection{Finite-shot identifiability}\label{finite-shot-identifiability}

The distinction of model parameters from finitely many random observations, conditional on a specified state class, sampling protocol, calibration model, and nuisance-parameter restrictions. Structural injectivity alone does not provide a finite-sample error or confidence guarantee.

\nopagebreak[4]

\emph{Reference:} Calibrated Dynamical Identification in Fermionic Systems, Simultaneous confidence with a nuisance budget.

\subsection{Fixed-index Jacobi upper spectral edge}\label{fixed-index-jacobi-upper-spectral-edge}

The regime in which the \(k\)th Jacobi eigenvalue counted from the upper edge is studied as \(N\to\infty\) while the nonnegative index \(k\) remains fixed. Constants may depend on \(k\); a statement in this regime does not automatically apply to \(k=k(N)\).

\nopagebreak[4]

\emph{Reference:} Fixed-Index Jacobi Spectral Expansion Through Inverse-Square Order, §1.

\subsection{Fixed-index Wilson-majorant weight}\label{fixed-index-wilson-majorant-weight}

The expectation \(q_{N,k}(\mu)=\langle v_{N,k},G_Nv_{N,k}\rangle\) at the \(k\)th upper-edge Jacobi mode, with \(k\) and scaled mass \(\mu\) fixed in the stated asymptotic. A uniform weighted spectral law requires additional information.

\nopagebreak[4]

\emph{Reference:} Fixed-Index Wilson Majorant Weights and Longitudinal Edge Contributions, §3.

\subsection{Formal modular period}\label{formal-modular-period}

A coefficientwise WKB contour integral with specified curve, branch, cycle, orientation and leading normalization. Exact-differential reductions identify its formal coefficients. Their spectral interpretation additionally requires a connection theorem and a remainder estimate; coefficient identities alone are insufficient.

\nopagebreak[4]

\emph{Reference:} Higher-Order Modular Periods for the Half-Filled Sine Kernel, The exact differential equation and the formal period.

\subsection{Fractional Sobolev boundary range}\label{fractional-sobolev-boundary-range}

The range \(0<s<1/2\) used for the interval-cutoff and nuclear output-regularity estimates in the Wilson boundary analysis. It is a strict range; the endpoint \(s=1/2\) requires a separate theorem.

\nopagebreak[4]

\emph{Reference:} Cross-Boundary Trace-Class Estimates for Wilson--Dirac Quasifree States, ``direct \(0<s<1/2\) Sobolev nuclear estimate''.

\subsection{Gauge invariance of a state}\label{gauge-invariance-of-a-state}

Invariance \(\omega\circ\gamma_\theta=\omega\) under a specified gauge automorphism action. In a finite density-matrix representation with \(U(\theta)=e^{i\theta Q}\), this means \(U(\theta)\rho U(\theta)^*=\rho\). Restricted measurements alone do not establish invariance of the underlying state.

\nopagebreak[4]

\emph{Reference:} Restricted Observables and CAR-State Distinguishability, §1.

\subsection{Global reference-modular linear envelope}\label{global-reference-modular-linear-envelope}

An inequality bounding the reference modular matrix or its discrepancy by a constant multiple of \(|J_N|\) in their common functional calculus, uniformly over all finite reference modes. Such a linear envelope permits growth with \(N\) and is different from an \(N\)-independent bound for \(\|E_N\|\).

\nopagebreak[4]

\emph{Reference:} Uniform Bulk Modular Asymptotics and the Leading Wilson Entropy Law, Global linear bounds.

\subsection{Half-filled sine-kernel reference}\label{half-filled-sine-kernel-reference}

The scalar matrix \(B_N\) with \(B_{pq}=\sin[\pi(p-q)/2]/[\pi(p-q)]\) for \(p\ne q\) and \(B_{pp}=1/2\). For the specified massless infinite-chain Wilson projector its finite-interval covariance compression is related to \(B_N\) by an explicit interval-preserving unitary. A commuting tridiagonal Jacobi operator enables reference-mode calculations.

\nopagebreak[4]

\emph{Reference:} Exact Reference Geometry and the Two-Bond Asymptotic of a Wilson--Dirac Covariance, Theorem 1 and Corollary 2.

\subsection{Hilbert--Schmidt norm}\label{hilbertschmidt-norm}

The Schatten-2 norm, \(\|A\|_2^2=\operatorname{Tr}(A^*A)\) where finite. It is distinct from the trace norm \(\|A\|_1\) and the operator norm \(\|A\|\).

\nopagebreak[4]

\emph{Reference:} A Hilbert--Schmidt Scaling Limit for Finite-Interval Wilson--Dirac Covariances, §1.

\subsection{Initial Wronskians for an extended-Chebyshev system}\label{initial-wronskians-for-an-extended-chebyshev-system}

For an ordered smooth family \(f_1,\ldots,f_n\), the determinants \(W(f_1,\ldots,f_k)\) for \(k=1,\ldots,n\) are the initial Wronskians. An extended-Chebyshev criterion must specify the relevant interval and regularity assumptions; a single full determinant is not the entire criterion.

\nopagebreak[4]

\emph{Reference:} Positive Spectral Measures Under Finite Moment and Laplace Constraints, §2.

\subsection{Integrated Riccati residual certificate}\label{integrated-riccati-residual-certificate}

For a logarithmic derivative \(w=hu_-^{\prime}/u_-\) and finite approximate momentum \(S_K\), set \(w_K=S_K-(h/2)(\log S_K)^{\prime}\) and \(R_K=hw_K^{\prime}+w_K^2-(Q_t+h^2v)\). The certificate consists of bounds on this residual together with the positivity and integrability hypotheses needed to control the exact solution. Formal coefficient cancellation alone is not such a certificate.

\nopagebreak[4]

\emph{Reference:} Uniform Bulk Modular Asymptotics and the Leading Wilson Entropy Law, Integrated stability bound.

\subsection{Inverse modulus}\label{inverse-modulus}

A quantitative inequality bounding parameter separation by differences of specified observable responses on a declared parameter domain; a nonsingular local Jacobian alone is not a global inverse modulus.

\nopagebreak[4]

\emph{Reference:} Local Number-Conserving Response and Inference, A stable short-time map in the squared scale.

\subsection{Jacobi spectral majorant}\label{jacobi-spectral-majorant}

A positive operator \(G_N\) that dominates specified signed diagonal expectations and induces the positive measure \(\nu_N(A)=\operatorname{Tr}[G_N\mathbf1_A(|J_N|)]\). Bounds for its total mass and spectral tail require the stated Wilson projector construction.

\nopagebreak[4]

\emph{Reference:} Exact Spectral Majorants for a Wilson-Dirac Interval, Theorems 1 and 3.

\subsection{Leading Wilson entropy law}\label{leading-wilson-entropy-law}

The leading fixed-scaled-mass asymptotic of \(\mathcal D(C_{L,\mu/L}\Vert C_{L,0})\), with the massive covariance in the first argument and \(\mu>0\) fixed before \(L\to\infty\). Its coefficient, remainder bound and hypotheses belong to the entropy theorem. A bounded remainder and a convergent finite part are distinct conclusions.

\nopagebreak[4]

\emph{Reference:} Relative Entropy of Finite-Interval Wilson--Dirac Quasifree States, Leading Wilson relative entropy.

\subsection{Longitudinal and transverse Wilson components}\label{longitudinal-and-transverse-wilson-components}

The model-specific splitting \(\Delta_N=S_N+D_N\) of the massive-minus-massless compressed covariance, obtained from the full-line projector geometry. The words longitudinal and transverse refer to that splitting, not to a generic commuting decomposition of arbitrary compressed covariances.

\nopagebreak[4]

\emph{Reference:} Exact Spectral Majorants for a Wilson-Dirac Interval, §2.

\subsection{Maximal compatible support threshold}\label{maximal-compatible-support-threshold}

For nonnegative normalized measures compatible with a specified finite set of moment and Laplace observations, the supremum of their lower support endpoints. It is an extremal feasibility quantity, not necessarily the actual lower support endpoint of the measure that generated the observations.

\nopagebreak[4]

\emph{Reference:} Positive Spectral Measures Under Finite Moment and Laplace Constraints, §1.

\subsection{Modular reference discrepancy}\label{modular-reference-discrepancy}

The commuting reference matrix \(E_N=K_N+N\eta(J_N/N)\), where \(K_N=\log(I-B_N)-\log B_N\). Thus \(-N\eta(J_N/N)\) is an approximation to, rather than the definition of, the exact finite modular matrix. Weighted trace bounds and operator-norm bounds are different claims.

\nopagebreak[4]

\emph{Reference:} Uniform Bulk Modular Asymptotics and the Leading Wilson Entropy Law, Statement and conventions.

\subsection{Moving-window ultraviolet limit}\label{moving-window-ultraviolet-limit}

An asymptotic comparison in which the observation parameter \(s_\gamma\) changes with regulator \(\gamma\) and the relevant representation labels also increase with \(\gamma\). Pointwise convergence at each fixed representation label does not justify a trace-ideal limit without uniform error and summable-tail control.

\nopagebreak[4]

\emph{Reference:} Compact-Group Transfer Spectra and Ultraviolet Operator Limits, §2.

\subsection{Normalized heat state}\label{normalized-heat-state}

\(\rho_A(s)=e^{-sA}/\operatorname{Tr}(e^{-sA})\) for positive s and finite partition trace. It is invariant under scalar shifts \(A\mapsto A+cI\), unlike the unnormalized heat operator; perturbation bounds for the two objects are not interchangeable.

\nopagebreak[4]

\emph{Reference:} Compact-Group Transfer Spectra and Ultraviolet Operator Limits, §5.

\subsection{Nuclear output regularity}\label{nuclear-output-regularity}

Trace-class mapping into a stronger output space, such as \(\mathcal S_1(L^2,H^s)\). The domain, codomain, embedding and norm are part of the statement. It is stronger information than merely treating the same operator as an \(L^2\)-to-\(L^2\) map.

\nopagebreak[4]

\emph{Reference:} Cross-Boundary Trace-Class Estimates for Wilson--Dirac Quasifree States, ``direct \(0<s<1/2\) Sobolev nuclear estimate''.

\subsection{Observation completeness}\label{observation-completeness}

Injectivity of the map from admissible states to the complete vector of specified observations; relative to a model class.

\nopagebreak[4]

\emph{Reference:} Restricted Observables and CAR-State Distinguishability, The exact descent criterion.

\subsection{Observational equivalence}\label{observational-equivalence}

Two preparations induce equal statistics for every admissible observable under an explicitly specified protocol.

\nopagebreak[4]

\emph{Reference:} Restricted Observables and CAR-State Distinguishability, Exact equivalence on the full neutral algebra.

\subsection{Occupation-weighted modular variance}\label{occupation-weighted-modular-variance}

The nonnegative trace \(\operatorname{Tr}[B_N(I-B_N)E_N^2]\) for the commuting reference matrices. A bound for this weighted quantity does not by itself bound \(\|E_N\|\) over the full spectrum.

\nopagebreak[4]

\emph{Reference:} Uniform Bulk Modular Asymptotics and the Leading Wilson Entropy Law, Positive weighted modular control.

\subsection{Open-boundary truncation}\label{open-boundary-truncation}

Compression or termination of a translation-invariant hopping Hamiltonian to a finite open device. It generally does not preserve every algebraic relation among infinite-volume spectral projectors.

\nopagebreak[4]

\emph{Reference:} Projector-Equivalent Wilson--Dirac Hamiltonians and Dynamical Identification, §5.

\subsection{Ordered discrete-prolate reference modes}\label{ordered-discrete-prolate-reference-modes}

The common eigenvectors of the concentration and Jacobi matrices, indexed in a declared spectral order. At half filling the normalization is \(B_N=p(N,1/4)\) and \(J_N=2\sigma(N,1/4)/N\) in the notation of Slepian. Ordering is a finite spectral statement, distinct from a uniform modular approximation.

\nopagebreak[4]

\emph{Reference:} Half-Filled Sine-Kernel Endpoint and Extreme Modular Discrepancy, §5.

\subsection{Particle--hole involution}\label{particlehole-involution}

The alternating-sign unitary \(\Gamma_N=\operatorname{diag}(1,-1,1,\ldots)\) used to relate the occupied and unoccupied reference sectors. The operator identities involving \(B_N\), \(J_N\), and the Wilson majorant are proved in the cited source; a cancellation for the majorant does not identify a cancellation for the longitudinal perturbation.

\nopagebreak[4]

\emph{Reference:} Exact Reference Occupation Weights and Particle-Hole Pairing, §1.

\subsection{Positive differential tail operator}\label{positive-differential-tail-operator}

The self-adjoint Darboux-deleted realization with differential expression \(T_N=-D_z^2+(4N^2-1)/(16p)+p/16\), where \(dz/dx=(\cos x)^{-1/2}\) and \(p(z)=\cos x(z)\). Its domain is part of the definition. Equality of its spectrum with the omitted original spectrum requires the deletion theorem.

\nopagebreak[4]

\emph{Reference:} Positive Moment Bounds and Growing Plunge Windows, Exact deletion of the finite block.

\subsection{Positive Jacobi-weighted moment}\label{positive-jacobi-weighted-moment}

The nonnegative trace \(Q_N=\operatorname{Tr}(G_N|J_N|)\). It is a moment of the positive Jacobi spectral measure induced by \(G_N\), not the signed first moment \(\operatorname{Tr}(S_NJ_N)\).

\nopagebreak[4]

\emph{Reference:} Positive Reference-Modular Localization and Bounds for the Wilson Majorant Jacobi Moment, §3.

\subsection{Positive reciprocal-trace representation}\label{positive-reciprocal-trace-representation}

A Laplace or Stieltjes representation of the first omitted reciprocal trace \(Z_N(1)=\sum_{k\ge N}\mu_{N,k}^{-1}\) by a positive kernel or measure. Positivity can support signed finite-order remainders; it does not mean that an infinite asymptotic series converges.

\nopagebreak[4]

\emph{Reference:} Reciprocal Traces and the Half-Filled Plunge Law, Positive Laplace formula; Positive Moment Bounds and Growing Plunge Windows, Stieltjes representation and signed remainder.

\subsection{Positive spectral deformation}\label{positive-spectral-deformation}

Multiplication of a gapped Bloch Hamiltonian by a strictly positive momentum-dependent scalar, leaving occupied and unoccupied eigenprojectors unchanged while generally altering dispersion.

\nopagebreak[4]

\emph{Reference:} Projector-Equivalent Wilson--Dirac Hamiltonians and Dynamical Identification, Theorem 1.

\subsection{Positive Wilson reference-modular trace}\label{positive-wilson-reference-modular-trace}

The nonnegative trace \(P_N=\operatorname{Tr}(G_N|K_N|)\) for the positive Wilson majorant \(G_N\) and reference modular matrix \(K_N\). Its relation to the signed entropy decomposition and its asymptotic coefficient are separate theorem statements.

\nopagebreak[4]

\emph{Reference:} A Positive Reference-Modular Trace for Wilson--Dirac Intervals, §1.

\subsection{Precision replication and rigorous numerical enclosure}\label{precision-replication-and-rigorous-numerical-enclosure}

Agreement under increased working precision supports numerical stability at the retained resolution. A rigorous enclosure additionally requires exact arithmetic or justified outward rounding, including all discretization, spectral and truncation errors relevant to the claimed quantity.

\nopagebreak[4]

\emph{Reference:} Positive Spectral Measures Under Finite Moment and Laplace Constraints, §3; Uniform Bulk Modular Asymptotics and the Leading Wilson Entropy Law, Numerical checks and remaining scope.

\subsection{Projector-equivalent Hamiltonians}\label{projector-equivalent-hamiltonians}

Hamiltonians whose specified spectral projections coincide under a common algebra identification. Their equal-time filled-band quasi-free states can agree while their real-time dynamical generators differ.

\nopagebreak[4]

\emph{Reference:} Projector-Equivalent Wilson--Dirac Hamiltonians and Dynamical Identification, §§1 and 5.

\subsection{Quasi-free state}\label{quasi-free-state}

A fermionic Gaussian CAR state whose field correlations obey the signed fermionic Wick pairing rule, equivalently Pfaffian form. For gauge-invariant quasi-free states, suitably ordered number-balanced creation/annihilation correlations reduce to determinants.

\nopagebreak[4]

\emph{Reference:} Operational Identification and Calibration in Two-Mode Quasifree Systems, §1.

\subsection{Reference occupation fluctuation}\label{reference-occupation-fluctuation}

For the half-filled sine-kernel covariance \(B_N\), the scalar \(\operatorname{Tr}[B_N(I-B_N)]\) is the total particle-number fluctuation in the gauge-invariant quasifree state. Its exact evaluation and large-\(N\) behavior are results, not part of this definition.

\nopagebreak[4]

\emph{Reference:} Exact Reference Occupation Weights and Particle-Hole Pairing, §2, Proposition.

\subsection{Reference-crossing integral envelope}\label{reference-crossing-integral-envelope}

The scalar integral \(I_N=\frac12\int_0^1\sqrt{1-u^N}/[\sqrt u(1-u)]\,du\) for positive even \(N\), used to bound the nuclear norm of the half-filled reference projection block crossing a finite interval boundary. Its estimates belong to the cited theorem; this integral is not an entropy finite part.

\nopagebreak[4]

\emph{Reference:} Finite-Interval Occupation Logarithm and Reference-Crossing Bounds, Theorem 2.

\subsection{Representation multiplicity in heat traces}\label{representation-multiplicity-in-heat-traces}

The multiplicity with which each spectral label contributes to a specified trace. For the central SU(2) model, multiplicities are 1 on the character subspace and \(n^2\) on full \(L^2(SU(2))\); changing this choice changes Schatten asymptotics.

\nopagebreak[4]

\emph{Reference:} Compact-Group Transfer Spectra and Ultraviolet Operator Limits, §§1 and 2.

\subsection{Schatten-1 cross-boundary control}\label{schatten-1-cross-boundary-control}

An estimate in trace norm for a block \(\chi(P_m-P_0)(I-\chi)\) mapping the outside of an interval into the inside. Such an estimate can enter entropy continuity arguments, but its norm, model, mass scaling and boundary channels must all be specified.

\nopagebreak[4]

\emph{Reference:} Uniform Square-Root-Mass Control of Wilson-Dirac Boundary Covariances, Theorem 1.

\subsection{Selected and omitted differential spectrum}\label{selected-and-omitted-differential-spectrum}

The split at index \(N\) in the ordered spectrum of the Friedrichs realization of \(-(\cos x\,u^{\prime})^{\prime}-[(N^2-1)/4]\cos x\,u\) on \((-\pi/2,\pi/2)\). The omitted eigenvalues are written \(\mu_{N,k}\), \(k\ge N\), and their reciprocal traces are \(Z_N(s)=\sum_{k\ge N}\mu_{N,k}^{-s}\) for \(s>1/2\). Positivity and convergence, and identification of the selected block with the finite Jacobi modes, are results in the cited paper.

\nopagebreak[4]

\emph{Reference:} Spectral Tails and Exact Concentration Logits, The differential spectrum and its first finite block.

\subsection{Sharp occupation-logarithm inequality}\label{sharp-occupation-logarithm-inequality}

A universal scalar inequality comparing \(b(1-b)\log^2[(1-b)/b]\) with \(\sqrt{b(1-b)}\) for \(0<b<1\). Its optimal coefficient is denoted \(C_{\log}\); the cited theorem expresses it using the positive root \(x_*\) of \(x_*\tanh x_*=2\). The optimization domain is the entire interval \((0,1)\) unless a smaller occupation class is explicitly specified.

\nopagebreak[4]

\emph{Reference:} Finite-Interval Occupation Logarithm and Reference-Crossing Bounds, Theorem 1.

\subsection{Signed half-filled plunge law}\label{signed-half-filled-plunge-law}

The central-mode comparison expressed in the sign convention \(e\sim-\varepsilon/(4N^2)\), equivalently \(e\sim\pi\lambda/(4N^2)\) in the regime specified by the cited theorem. A use of this named asymptotic must specify the uniform eigenvalue window and relative-error estimate from the corresponding theorem; it is not a general-filling convention.

\nopagebreak[4]

\emph{Reference:} Positive Moment Bounds and Growing Plunge Windows, Uniform center law beyond shrinking eigenvalues.

\subsection{Simultaneous confidence region}\label{simultaneous-confidence-region}

A random set covering the true parameter or prediction vector with the specified joint probability under sampling and calibration assumptions; separate marginal confidence levels cannot simply be reused as a common joint level.

\nopagebreak[4]

\emph{Reference:} Calibrated Dynamical Identification in Fermionic Systems, Simultaneous confidence with a nuisance budget.

\subsection{Spectral convergence}\label{spectral-convergence}

Specified eigenvalue/operator-spectrum convergence with topology and boundary assumptions.

\nopagebreak[4]

\emph{Reference:} Spectral Tails and Exact Concentration Logits, Compact spectrum and elementary bounds.

\subsection{Spectral-wing localization}\label{spectral-wing-localization}

A decomposition of a reference spectral trace into an interior window and its complementary spectral wing, with an estimate on the omitted contribution. A moving cutoff must be stated as a function of \(N\) together with its quantifiers; fixed-index information is not a substitute.

\nopagebreak[4]

\emph{Reference:} Positive Reference-Modular Localization and Bounds for the Wilson Majorant Jacobi Moment, §2.

\subsection{Strong operator topology}\label{strong-operator-topology}

Pointwise norm convergence on each vector, weaker than operator-norm convergence.

\nopagebreak[4]

\emph{Reference:} Compact-Group Transfer Spectra and Ultraviolet Operator Limits, §8.

\subsection{Structural identifiability}\label{structural-identifiability}

Injectivity of a specified noiseless parameter-to-observation mapping, after declaring admissible parameter sets and equivalences. A locally nonsingular Jacobian does not exclude global aliases or establish completeness of the assumed model family.

\nopagebreak[4]

\emph{Reference:} Projector-Equivalent Wilson--Dirac Hamiltonians and Dynamical Identification, §4.

\subsection{Trace norm versus operational trace distance}\label{trace-norm-versus-operational-trace-distance}

For density operators \(\rho\) and \(\sigma\), the unhalved Schatten trace norm is \(\|\rho-\sigma\|_1=\operatorname{Tr}|\rho-\sigma|\). Conventional operational trace distance is \(\frac12\|\rho-\sigma\|_1\). When comparing state-distinguishability bounds one must check which normalization is being used.

\nopagebreak[4]

\emph{Reference:} Operational Identification and Calibration in Two-Mode Quasifree Systems, §1.

\subsection{Two-bond trace contribution}\label{two-bond-trace-contribution}

The linear branch \(\mathcal P_L=\operatorname{Tr}[(A_N-B_N)(-\pi J_N)]\) in the Wilson finite relative-entropy decomposition. The trace reduces exactly to intra-cell and inter-cell nearest-neighbor Fourier coefficients; its fixed-positive-mass asymptotic follows from complete elliptic integrals.

\nopagebreak[4]

\emph{Reference:} Exact Reference Geometry and the Two-Bond Asymptotic of a Wilson--Dirac Covariance, Theorem 3.

\subsection{Uniform compact-interior modular expansion}\label{uniform-compact-interior-modular-expansion}

A finite-order spectral expansion whose error bound is uniform for actual modes with \(|\lambda/N|\le\tau\), where \(0<\tau<1/2\) and the truncation order are fixed before \(N\to\infty\). Uniformity in a moving edge window or a growing truncation order is a stronger requirement.

\nopagebreak[4]

\emph{Reference:} Uniform Bulk Modular Asymptotics and the Leading Wilson Entropy Law, Uniform interior expansion.

\subsection{Validated local root enclosure}\label{validated-local-root-enclosure}

A specified interval box together with justified outward-rounded arithmetic and a contraction or interval-Newton/Krawczyk condition establishing a root, and sometimes uniqueness, in that box. Local enclosure is not global optimality; any global result requires an additional analytical argument.

\nopagebreak[4]

\emph{Reference:} Positive Spectral Measures Under Finite Moment and Laplace Constraints, §3.

\subsection{Whole-line weighted resolvent}\label{whole-line-weighted-resolvent}

The form/resolvent transformation with \(p=\operatorname{sech}r\), \(h=1/N\) and \(P_h(t)=-h^2D_r^2+(1-2t\operatorname{sech}r)/4+h^2\operatorname{sech}^2r/4\). Multiplication by \(p^{1/2}\) appears on both sides of the ordinary whole-line inverse. This is not an identification of its ordinary spectrum with the discrete tail spectrum.

\nopagebreak[4]

\emph{Reference:} Uniform Bulk Modular Asymptotics and the Leading Wilson Entropy Law, Weighted resolvent identity.

\subsection{Wick identity}\label{wick-identity}

The signed sum over pairings that expresses even quasi-free CAR field correlations in terms of two-point functions, equivalently a Pfaffian. Determinant form requires the gauge-invariant, suitably ordered number-balanced specialization; it is not the general CAR rule.

\nopagebreak[4]

\emph{Reference:} Operational Identification and Calibration in Two-Mode Quasifree Systems, §1.

\subsection{Acknowledgments and contributions}\label{acknowledgments-and-contributions}

Kirtus Justus provided the research direction and publication requirements. ChatGPT (OpenAI) assisted with mathematical derivations, verification and editorial preparation, including the Wilson continuation. Claude (Anthropic), Gemini (Google), Grok (SpaceXAI), and DeepSeek (High-Flyer) are acknowledged for participation in the wider research program, without unsupported attribution of particular proofs. The named human author assumes responsibility for the scholarly content. These acknowledgments imply neither provider endorsement nor independent peer review.

\end{document}
